\documentclass[12pt]{article}
\usepackage[utf8]{inputenc}
\usepackage[T1]{fontenc}
\usepackage[spanish]{babel}
\usepackage{geometry}
\usepackage{amsmath,amssymb,amsthm,mathtools}
\usepackage{enumitem}

\geometry{margin=2.5cm}

\theoremstyle{definition}
\newtheorem{definition}{Definición}[section]
\newtheorem{example}[definition]{Ejemplo}
\newtheorem{exercise}[definition]{Ejercicio}
\theoremstyle{plain}
\newtheorem{theorem}[definition]{Teorema}
\newtheorem{proposition}[definition]{Proposición}
\newtheorem{lemma}[definition]{Lema}
\newtheorem{corollary}[definition]{Corolario}
\theoremstyle{remark}
\newtheorem*{remark}{Observación}

\begin{document}

\begin{center}
{\Large Análisis Matemático I}\\[2mm]
{\large Apuntes organizados por fecha}\\[2mm]
\end{center}

\begin{remark}
La notación y la terminología se normalizan siguiendo \textit{Rudin, Principios de Análisis Matemático}, 3.\textsuperscript{a} edición, para mantener consistencia en los conceptos de conjuntos numerables, espacios métricos y topología básica.
\end{remark}

\section*{11May26}

\subsection*{Imágenes, preimágenes y equivalencia}

\begin{definition}[Imagen directa]
Sean $A$ y $B$ conjuntos, y sea $f\colon A\to B$.
Si $C\subset A$, la imagen directa de $C$ bajo $f$ es
\[
f(C)=\{f(a)\mid a\in C\}.
\]
\end{definition}

\begin{definition}[Imagen inversa]
Sean $A$ y $B$ conjuntos, y sea $f\colon A\to B$.
Si $D\subset B$, la imagen inversa de $D$ bajo $f$ es
\[
f^{-1}(D)=\{a\in A\mid f(a)\in D\}.
\]
\end{definition}

\begin{proposition}[Propiedades básicas]
Sea $f\colon A\to B$.
\begin{enumerate}[label=\arabic*.]
    \item $f\!\left(\bigcup_\alpha C_\alpha\right)=\bigcup_\alpha f(C_\alpha)$.
    \item $f^{-1}\!\left(\bigcup_\alpha D_\alpha\right)=\bigcup_\alpha f^{-1}(D_\alpha)$.
    \item $f^{-1}\!\left(\bigcap_\alpha D_\alpha\right)=\bigcap_\alpha f^{-1}(D_\alpha)$.
    \item $f^{-1}(D^c)=\bigl(f^{-1}(D)\bigr)^c$.
\end{enumerate}
\end{proposition}

\begin{example}
Para toda función $f$, se tiene
\[
f(\varnothing)=\varnothing,
\qquad
f^{-1}(\varnothing)=\varnothing.
\]
\end{example}

\begin{definition}[Conjuntos equivalentes]
Sean $A$ y $B$ conjuntos.
Decimos que $A$ y $B$ son equivalentes si existe una biyección $f\colon A\to B$.
\end{definition}

\begin{remark}
La equivalencia anterior define una relación de equivalencia entre conjuntos.
\end{remark}

\subsection*{Cardinalidad y numerabilidad}

\begin{definition}[Cardinalidad]
Sea $A$ un conjunto. Su cardinalidad es la clase de equivalencia a la que pertenece.
\end{definition}

\begin{definition}[Conjunto finito]
Un conjunto $A$ es finito si existe $n\in\mathbb N$ tal que $A$ es biyectivo con $\{1,\dots,n\}$.
\end{definition}

\begin{definition}[Conjunto numerable]
Un conjunto $A$ es numerable si existe una biyección $f\colon\mathbb N\to A$.
\end{definition}

\begin{theorem}[Subconjuntos infinitos de conjuntos numerables]
Si $H$ es un conjunto numerable y $E\subset H$ es infinito, entonces $E$ es numerable.
\end{theorem}

\begin{proof}
Sea $f\colon\mathbb N\to H$ una biyección.
Como $E$ es infinito, podemos construir por inducción una sucesión estrictamente creciente
\[
h(1)<h(2)<h(3)<\cdots
\]
tal que $f(h(j))\in E$ para todo $j\in\mathbb N$.
Entonces $f\circ h\colon\mathbb N\to E$ es biyectiva.
\end{proof}

\begin{lemma}
$\mathbb N\times\mathbb N$ es numerable.
\end{lemma}

\begin{proof}
Se enumera por diagonales:
\[
(1,1),(1,2),(2,1),(1,3),(2,2),(3,1),\dots
\]
lo cual induce una biyección con $\mathbb N$.
\end{proof}

\begin{theorem}[Unión numerable de numerables]
Si $\{E_n\}_{n=1}^\infty$ es una familia de conjuntos numerables, entonces
\[
\bigcup_{n=1}^\infty E_n
\]
es numerable.
\end{theorem}

\begin{proof}
Primero, si los $E_n$ son mutuamente ajenos, definimos
\[
T\colon\mathbb N\times\mathbb N\to \bigcup_{n=1}^\infty E_n,
\qquad T(n,j)=x_{nj},
\]
donde $E_n=\{x_{n1},x_{n2},\dots\}$.
El mapa $T$ es biyectivo.
En el caso general, se reemplaza $E_n$ por
\[
F_1=E_1,\qquad F_n=E_n\setminus\bigcup_{k=1}^{n-1}E_k,\quad n\ge2,
\]
y se aplica el caso disjunto.
\end{proof}

\begin{corollary}
Si $A$ es numerable y, para cada $\alpha\in A$, el conjunto $E_\alpha$ es numerable, entonces
\[
\bigcup_{\alpha\in A}E_\alpha
\]
es numerable.
\end{corollary}

\begin{corollary}[Producto cartesiano de numerables]
Si $A$ y $B$ son conjuntos numerables, entonces $A\times B$ es numerable.
\end{corollary}

\begin{proof}
Para cada $\alpha\in A$, sea
\[
A_\alpha=\{(\alpha,\beta)\mid \beta\in B\}.
\]
Entonces
\[
A\times B=\bigcup_{\alpha\in A}A_\alpha,
\]
y cada $A_\alpha$ es numerable porque está en biyección con $B$.
\end{proof}

\begin{theorem}[Los racionales son numerables]
\[
\mathbb Q=\left\{\frac{a}{b}\,\middle|\, a\in\mathbb Z,\ b\in\mathbb N,\ b>0\right\}
\]
es numerable.
\end{theorem}

\begin{proof}
La aplicación
\[
S\colon\mathbb Z\times\mathbb N\to\mathbb Q,\qquad S(a,b)=\frac ab,
\]
es sobreyectiva.
Como $\mathbb Z\times\mathbb N$ es numerable, también lo es su imagen $\mathbb Q$.
\end{proof}

\begin{corollary}
Si $A$ es numerable y $f\colon A\to B$ es sobreyectiva, entonces $B$ es a lo más numerable.
\end{corollary}

\begin{theorem}[Ejemplo no numerable]
Sea
\[
\mathcal C=\{f\colon\mathbb N\to\{0,1\}\}.
\]
Entonces $\mathcal C$ no es numerable.
\end{theorem}

\begin{proof}
Supongamos que $\mathcal C=\{f_1,f_2,\dots\}$.
Definimos
\[
g(n)=1-f_n(n).
\]
Entonces $g\in\mathcal C$, pero $g\neq f_n$ para todo $n\in\mathbb N$, contradicción.
\end{proof}

\subsection*{Espacios métricos y espacios de funciones}

\begin{definition}[Métrica]
Sea $X$ un conjunto no vacío.
Una métrica en $X$ es una función
\[
d\colon X\times X\to [0,\infty)
\]
tal que, para todo $p,q,r\in X$:
\begin{enumerate}[label=\roman*)]
    \item $d(p,q)\ge0$ y $d(p,q)=0$ si y solo si $p=q$;
    \item $d(p,q)=d(q,p)$;
    \item $d(p,r)\le d(p,q)+d(q,r)$.
\end{enumerate}
\end{definition}

\begin{example}[Distancia en $\mathbb R$]
En $\mathbb R$, la función $d(a,b)=|a-b|$ es una métrica.
\end{example}

\begin{example}[Distancia euclídea]
En $\mathbb R^d$,
\[
d(a,b)=\sqrt{\sum_{i=1}^d(a_i-b_i)^2}
\]
es una métrica.
\end{example}

\begin{example}[Métrica discreta]
Si $X$ es cualquier conjunto, la función
\[
d(x,y)=
\begin{cases}
1,& x\neq y,\\
0,& x=y,
\end{cases}
\]
define una métrica.
\end{example}

\begin{definition}[Funciones continuas]
Sean $a,b\in\mathbb R$ con $a<b$.
Denotamos por
\[
C[a,b]=\{f:[a,b]\to\mathbb R\mid f \text{ es continua}\}.
\]
\end{definition}

\begin{theorem}[Métrica del supremo en funciones continuas]
Sea
\[
d\colon C[a,b]\times C[a,b]\to[0,\infty),\qquad
d(f,g)=\sup\{|f(x)-g(x)|\mid x\in[a,b]\}.
\]
Entonces $d$ es una métrica.
\end{theorem}

\begin{proof}
La positividad y la simetría son inmediatas.
Si $f\neq g$, existe $x_0\in[a,b]$ tal que $f(x_0)\neq g(x_0)$; por continuidad, eso produce un intervalo $U$ donde $f\neq g$ y por tanto $d(f,g)>0$.
La desigualdad del triángulo se obtiene de
\[
|f(x)-g(x)|\le |f(x)-h(x)|+|h(x)-g(x)|
\]
y al tomar supremo en $x$.
\end{proof}

\begin{definition}[Funciones acotadas]
Sea $(X,d)$ un espacio métrico y sea $A$ un conjunto cualquiera.
Denotamos por
\[
\mathcal B(A,X)=\{f\colon A\to X\mid f \text{ es acotada}\}.
\]
\end{definition}

\begin{definition}[Métrica del supremo en funciones acotadas]
Si $f,g\in\mathcal B(A,X)$, definimos
\[
D(f,g)=\sup\{d(f(x),g(x))\mid x\in A\}.
\]
\end{definition}

\begin{theorem}
La función $D$ es una métrica en $\mathcal B(A,X)$.
\end{theorem}

\begin{proof}
La positividad y la simetría son inmediatas.
Para la desigualdad del triángulo, basta usar la desigualdad del triángulo en $(X,d)$ y tomar supremo.
\end{proof}

\begin{example}
Si $A=\{1,2\}$ y $X=\mathbb R$, entonces $\mathcal B(A,X)\cong\mathbb R^2$.
Si identificamos $f$ con $(f(1),f(2))$, la métrica $D$ se escribe como
\[
D(f,g)=\max\{|f(1)-g(1)|,\ |f(2)-g(2)|\}.
\]
\end{example}

\section*{12May26}

\subsection*{Espacios de funciones y métrica del supremo}

\begin{theorem}[Métrica del supremo en funciones continuas]
Sean $a,b\in\mathbb R$ con $a<b$ y sea
\[
d(f,g)=\sup\{\,|f(x)-g(x)|\mid x\in[a,b]\,\}.
\]
Entonces $d$ es una métrica en $C[a,b]$.
\end{theorem}

\begin{proof}
La positividad y la simetría son inmediatas. Si $f\neq g$, existe $x_0\in[a,b]$ tal que $f(x_0)\neq g(x_0)$; en particular,
\[
d(f,g)\ge |f(x_0)-g(x_0)|>0.
\]
La desigualdad triangular se obtiene de
\[
|f(x)-g(x)|\le |f(x)-h(x)|+|h(x)-g(x)|
\]
y al tomar supremo en $x\in[a,b]$.
\end{proof}

\begin{definition}[Funciones acotadas]
Sea $(X,d)$ un espacio métrico y sea $A$ un conjunto cualquiera. Denotamos por
\[
\mathcal B(A,X)=\{f\colon A\to X\mid f \text{ es función acotada}\}.
\]
\end{definition}

\begin{definition}[Métrica del supremo en funciones acotadas]
Si $f,g\in\mathcal B(A,X)$, definimos
\[
D(f,g)=\sup\{d(f(x),g(x))\mid x\in A\}.
\]
\end{definition}

\begin{theorem}
La función $D$ es una métrica en $\mathcal B(A,X)$.
\end{theorem}

\begin{proof}
La positividad y la simetría son inmediatas. Para la desigualdad del triángulo, usamos la desigualdad triangular en $(X,d)$ y tomamos supremo.
\end{proof}

\begin{example}
Si $A=\{1,2\}$ y $X=\mathbb R$, entonces $\mathcal B(A,X)\cong\mathbb R^2$. Si identificamos $f$ con $(f(1),f(2))$, la métrica $D$ se escribe como
\[
D(f,g)=\max\{|f(1)-g(1)|,\ |f(2)-g(2)|\}.
\]
\end{example}

\subsection*{Vecindades, puntos límite e interior}

\begin{definition}[Vecindad]
Sea $(X,d)$ un espacio métrico, sea $p\in X$ y sea $r>0$.
La vecindad o entorno de radio $r$ centrada en $p$ es
\[
N_r(p)=\{q\in X\mid d(q,p)<r\}.
\]
\end{definition}

\begin{definition}[Punto límite]
Sea $E\subset X$ y sea $p\in X$.
Decimos que $p$ es punto límite de $E$ si toda vecindad de $p$ contiene elementos de $E$ distintos de $p$.
Equivalentemente,
\[
\forall r>0,\qquad N_r(p)\cap(E\setminus\{p\})\neq\varnothing.
\]
\end{definition}

\begin{definition}[Punto aislado]
Sea $E\subset X$ y sea $p\in E$.
Si $p$ no es punto límite de $E$, decimos que $p$ es un punto aislado de $E$.
\end{definition}

\begin{example}
En $X=\mathbb R$ y
\[
E=(0,1]\cup\{2\},
\]
los puntos $0$ y $1$ son puntos límite de $E$, mientras que $2$ es un punto aislado de $E$.
\end{example}

\begin{definition}[Punto interior]
Sea $E\subset X$ y sea $p\in X$.
Decimos que $p$ es un punto interior de $E$ si existe $r>0$ tal que
\[
N_r(p)\subset E.
\]
\end{definition}

\subsection*{Conjuntos abiertos, cerrados, perfectos, acotados y densos}

\begin{definition}[Conjunto abierto]
Un conjunto $E\subset X$ es abierto si todos sus elementos son puntos interiores.
\end{definition}

\begin{definition}[Conjunto cerrado]
Un conjunto $E\subset X$ es cerrado si contiene a todos sus puntos límite.
\end{definition}

\begin{definition}[Complemento]
Si $E\subset X$, su complemento es
\[
E^c=\{x\in X\mid x\notin E\}.
\]
\end{definition}

\begin{definition}[Conjunto perfecto]
Un conjunto $E\subset X$ es perfecto si es cerrado y todos sus elementos son puntos límite de $E$.
\end{definition}

\begin{definition}[Conjunto acotado]
Un conjunto $E\subset X$ es acotado si existen $M>0$ y $q\in X$ tales que
\[
d(p,q)<M\qquad \forall p\in E.
\]
\end{definition}

\begin{definition}[Conjunto denso]
Un conjunto $E\subset X$ es denso en $X$ si todo $p\in X$ es punto límite de $E$.
\end{definition}

\begin{example}
En $X=\mathbb R$, el conjunto
\[
E=(0,1]\cup\{2\}
\]
tiene a $2$ como punto aislado, a $0$ y $1$ como puntos límite, y no es perfecto.
\end{example}

\section*{14May26}

\subsection*{Vecindades y conjuntos derivados}

\begin{theorem}[Toda vecindad es abierta]
Sea $(X,d)$ un espacio métrico. Toda vecindad $N_r(p)$ es un conjunto abierto.
\end{theorem}

\begin{proof}
Sea $q\in N_r(p)$.
Tomamos
\[
h=r-d(q,p)>0.
\]
Si $a\in N_h(q)$, entonces
\[
d(a,p)\le d(a,q)+d(q,p)<h+d(q,p)=r.
\]
Por tanto, $a\in N_r(p)$, y así $N_h(q)\subset N_r(p)$.
\end{proof}

\begin{definition}[Conjunto derivado]
Sea $E\subset X$.
Denotamos por $E'$ al conjunto de todos los puntos límite de $E$.
\end{definition}

\begin{theorem}
Si $p\in E'$, entonces $N_r(p)\cap E$ es un conjunto infinito para todo $r>0$.
\end{theorem}

\begin{proof}
Supongamos lo contrario.
Entonces existe $r>0$ tal que $N_r(p)\cap E$ es finito.
Escribamos
\[
N_r(p)\cap E=\{q_1,\dots,q_n\}.
\]
Tomemos
\[
t=\min\{d(q_1,p),\dots,d(q_n,p)\}>0.
\]
Entonces
\[
N_t(p)\cap(E\setminus\{p\})=\varnothing,
\]
contradicción con que $p\in E'$.
\end{proof}

\begin{corollary}
Si $E\subset X$ es finito, entonces $E'=\varnothing$.
\end{corollary}

\begin{corollary}
Si $E\subset X$ es finito, entonces $E$ es cerrado.
\end{corollary}

\subsection*{Abiertos y cerrados}

\begin{theorem}
Sea $(X,d)$ un espacio métrico y sea $E\subset X$. Entonces $E$ es abierto si y solo si $E^c$ es cerrado.
\end{theorem}

\begin{proof}
Supongamos que $E^c$ es cerrado.
Sea $x\in E$.
Si $x$ no fuese interior, entonces para todo $r>0$ tendríamos
\[
N_r(x)\cap E^c\neq\varnothing,
\]
y así $x\in (E^c)'$, contradicción.
Luego existe $r>0$ tal que $N_r(x)\subset E$, y $E$ es abierto.

Recíprocamente, si $E$ es abierto, entonces $E^c$ es cerrado porque $(E^c)' \subset E^c$.
\end{proof}

\begin{theorem}
Si $G_1,\dots,G_n\subset X$ son abiertos, entonces
\[
G_1\cap\cdots\cap G_n
\]
es abierto.
\end{theorem}

\begin{proof}
Sea $x\in G_1\cap\cdots\cap G_n$.
Para cada $i$ existe $r_i>0$ tal que $N_{r_i}(x)\subset G_i$.
Tomamos $r=\min\{r_1,\dots,r_n\}$.
Entonces
\[
N_r(x)\subset G_1\cap\cdots\cap G_n.
\]
\end{proof}

\begin{theorem}
Si $F_1,\dots,F_n\subset X$ son cerrados, entonces
\[
F_1\cup\cdots\cup F_n
\]
es cerrado.
\end{theorem}

\begin{proof}
Como
\[
(F_1\cup\cdots\cup F_n)^c=F_1^c\cap\cdots\cap F_n^c,
\]
y cada $F_i^c$ es abierto, el resultado sigue del teorema anterior.
\end{proof}

\begin{proposition}[Operaciones arbitrarias]
\begin{enumerate}[label=\arabic*.]
    \item La unión arbitraria de abiertos es abierta.
    \item La intersección arbitraria de cerrados es cerrada.
\end{enumerate}
\end{proposition}

\begin{proof}
Se usan los complementos y el teorema abierto-cerrado.
\end{proof}

\begin{example}
Si
\[
E=\left\{\frac1n\mid n\in\mathbb N\right\},\quad
F=\mathbb R^2,\quad
G=(a,b),\quad
H=[a,b],\quad
I=[a,b),\quad
J=(a,b],
\]
entonces, según la tabla del pizarrón,
\[
\begin{array}{c|cccc}
 & \text{cerrado} & \text{abierto} & \text{perfecto} & \text{acotado}\\
\hline
E & \text{NO} & \text{NO} & \text{NO} & \text{SÍ}\\
F & \text{SÍ} & \text{SÍ} & \text{SÍ} & \text{NO}\\
G & \text{NO} & \text{SÍ} & \text{NO} & \text{SÍ}\\
H & \text{SÍ} & \text{NO} & \text{SÍ} & \text{SÍ}\\
I & \text{NO} & \text{NO} & \text{NO} & \text{SÍ}\\
J & \text{NO} & \text{NO} & \text{NO} & \text{SÍ}
\end{array}
\]
\end{example}

\section*{18May26}

\subsection*{Uniones e intersecciones}

\begin{theorem}[Intersección finita y unión finita]
Sea $(X,d)$ un espacio métrico.
\begin{enumerate}[label=\alph*.]
    \item La intersección finita de conjuntos abiertos es abierta.
    \item La unión finita de conjuntos cerrados es cerrada.
\end{enumerate}
\end{theorem}

\begin{proof}
El inciso (a) se demuestra tomando, para cada punto de la intersección, un radio mínimo entre los radios que lo mantienen dentro de cada abierto.
El inciso (b) se deduce del inciso (a) tomando complementos.
\end{proof}

\begin{example}
Para cada $n\in\mathbb N$, sea
\[
G_n=\left(-\frac1n,\frac1n\right).
\]
Entonces cada $G_n$ es abierto en $\mathbb R$ y
\[
\bigcap_{n=1}^\infty G_n=\{0\},
\]
que no es abierto.

Además,
\[
G_n^c=\left(-\infty,-\frac1n\right]\cup\left[\frac1n,\infty\right),
\qquad
\bigcup_{n=1}^\infty G_n^c=(-\infty,0)\cup(0,\infty),
\]
que es abierto pero no cerrado.
\end{example}

\subsection*{Interior, cerradura y derivado}

\begin{definition}[Conjunto derivado]
Sea $E\subset X$. Denotamos por $E'$ al conjunto de todos los puntos límite de $E$.
\end{definition}

\begin{definition}[Interior y cerradura]
Sea $E\subset X$. El interior de $E$ es
\[
E^\circ=\{x\in X\mid \exists r>0 \text{ tal que } N_r(x)\subset E\},
\]
y la cerradura de $E$ es
\[
\overline E=E\cup E'.
\]
\end{definition}

\begin{proposition}[Caracterización del interior]
El interior $E^\circ$ es el mayor conjunto abierto contenido en $E$.
Equivalentemente,
\[
E^\circ=\bigcup\{A\subset E\mid A \text{ es abierto}\}.
\]
\end{proposition}

\begin{proposition}[Caracterización de la cerradura]
La cerradura $\overline E$ es el menor conjunto cerrado que contiene a $E$.
\end{proposition}

\begin{example}
Sea
\[
E=\{(x,y)\in\mathbb R^2\mid x^2+y^2\le 1,\ y>0\}.
\]
Entonces $E$ no es cerrado ni abierto,
\[
E^\circ=\{(x,y)\in\mathbb R^2\mid x^2+y^2<1,\ y>0\},
\]
y
\[
\overline E=\{(x,y)\in\mathbb R^2\mid x^2+y^2\le 1,\ y\ge 0\}.
\]
\end{example}

\subsection*{Subespacios}

\begin{theorem}[Abiertos relativos]
Sean $(X,d)$ un espacio métrico y $Y\subset X$. Si $E\subset Y$, entonces $E$ es abierto relativo en $Y$ si y solo si existe un abierto $G\subset X$ tal que
\[
E=Y\cap G.
\]
\end{theorem}

\begin{proof}
Si $E$ es abierto relativo en $Y$, entonces para cada $x\in E$ existe $r_x>0$ tal que
\[
N_{r_x}(x)\cap Y\subset E.
\]
Por tanto
\[
E=\bigcup_{x\in E}\bigl(N_{r_x}(x)\cap Y\bigr)
=Y\cap\left(\bigcup_{x\in E}N_{r_x}(x)\right),
\]
y la unión de bolas abiertas es abierta en $X$.
Recíprocamente, si $E=Y\cap G$ con $G$ abierto en $X$, entonces para cada $x\in E$ existe $r>0$ tal que $N_r(x)\subset G$, y así $N_r(x)\cap Y\subset E$.
\end{proof}

\subsection*{Compacidad}

\begin{definition}[Compacidad]
Sea $(X,d)$ un espacio métrico y sea $K\subset X$.
Decimos que $K$ es compacto si toda cubierta abierta de $K$ admite una subcubierta finita.
\end{definition}

\begin{theorem}[Compacidad y subespacios]
Sean $(X,d)$ un espacio métrico, $Y\subset X$ y $K\subset Y$.
Entonces $K$ es compacto en $Y$ si y solo si $K$ es compacto en $X$.
\end{theorem}

\begin{proof}
Si $K$ es compacto en $Y$ y $\{W_\alpha\}$ es una cubierta abierta de $K$ en $X$, entonces $\{W_\alpha\cap Y\}$ es una cubierta abierta relativa de $K$ en $Y$.
Por compacidad de $K$ en $Y$, existe una subcubierta finita, y por tanto también una subcubierta finita de $K$ en $X$.

La recíproca se prueba tomando una cubierta abierta relativa de $K$ en $Y$, escribiendo cada abierto relativo como intersección con un abierto de $X$, y aplicando la compacidad de $K$ en $X$.
\end{proof}

\begin{theorem}[Los compactos son cerrados]
Si $(X,d)$ es un espacio métrico y $K\subset X$ es compacto, entonces $K$ es cerrado.
\end{theorem}

\begin{proof}
Probamos que $K^c$ es abierto.
Sea $p\in K^c$.
Para cada $q\in K$, definimos
\[
r_q=\frac12 d(p,q)>0.
\]
Entonces
\[
N_{r_q}(q)\cap N_{r_q}(p)=\varnothing.
\]
La familia $\{N_{r_q}(q)\}_{q\in K}$ cubre a $K$.
Como $K$ es compacto, existe una subcubierta finita
\[
K\subset N_{r_{q_1}}(q_1)\cup\cdots\cup N_{r_{q_n}}(q_n).
\]
Tomamos
\[
r=\min\{r_{q_1},\dots,r_{q_n}\}.
\]
Entonces $N_r(p)\cap K=\varnothing$, y por tanto $K^c$ es abierto.
\end{proof}

\section*{19May26}

\subsection*{Subconjuntos cerrados de compactos}

\begin{theorem}[Los compactos son cerrados]
Si $(X,d)$ es un espacio métrico y $K\subset X$ es compacto, entonces $K$ es cerrado.
\end{theorem}

\begin{proof}
Probamos que $K^c$ es abierto.
Sea $p\in K^c$. Para cada $q\in K$, tomamos
\[
r_q=\frac12 d(p,q)>0.
\]
Entonces
\[
N_{r_q}(q)\cap N_{r_q}(p)=\varnothing.
\]
La familia $\{N_{r_q}(q)\}_{q\in K}$ cubre a $K$.
Como $K$ es compacto, existe una subcubierta finita
\[
K\subset N_{r_{q_1}}(q_1)\cup\cdots\cup N_{r_{q_n}}(q_n).
\]
Tomando
\[
r=\min\{r_{q_1},\dots,r_{q_n}\},
\]
se obtiene $N_r(p)\cap K=\varnothing$.
Luego $K^c$ es abierto.
\end{proof}

\begin{theorem}[Subconjuntos cerrados de compactos]
Si $K\subset X$ es compacto y $F\subset K$ es cerrado, entonces $F$ es compacto.
\end{theorem}

\begin{proof}
Sea $\{V_\alpha\}_{\alpha\in\Lambda}$ una cubierta abierta de $F$ por abiertos de $X$.
Como $F$ es cerrado, $F^c$ es abierto.
Entonces
\[
\{V_\alpha\}_{\alpha\in\Lambda}\cup\{F^c\}
\]
es una cubierta abierta de $K$.
Por compacidad de $K$, existen $\alpha_1,\dots,\alpha_n\in\Lambda$ tales que
\[
K\subset V_{\alpha_1}\cup\cdots\cup V_{\alpha_n}\cup F^c.
\]
Como $F\subset K$, se sigue que
\[
F\subset V_{\alpha_1}\cup\cdots\cup V_{\alpha_n}.
\]
Por tanto, $F$ es compacto.
\end{proof}

\begin{corollary}
Si $F\subset X$ es cerrado y $K\subset X$ es compacto, entonces $F\cap K$ es compacto.
\end{corollary}

\begin{proof}
El conjunto $F\cap K$ es cerrado en $K$ y está contenido en el compacto $K$.
Por el teorema anterior, $F\cap K$ es compacto.
\end{proof}

\subsection*{Intersección finita}

\begin{theorem}[Propiedad de la intersección finita]
Sea $\{K_\alpha\}_{\alpha\in\Lambda}$ una familia de subconjuntos compactos de $X$.
Si toda subfamilia finita de $\{K_\alpha\}$ tiene intersección no vacía, entonces
\[
\bigcap_{\alpha\in\Lambda}K_\alpha\neq\varnothing.
\]
\end{theorem}

\begin{proof}
Supongamos que
\[
\bigcap_{\alpha\in\Lambda}K_\alpha=\varnothing.
\]
Fijemos $\alpha_0\in\Lambda$.
Entonces
\[
K_{\alpha_0}\cap\bigcap_{\alpha\in\Lambda\setminus\{\alpha_0\}}K_\alpha=\varnothing,
\]
de donde
\[
K_{\alpha_0}\subset\bigcup_{\alpha\in\Lambda\setminus\{\alpha_0\}}K_\alpha^c.
\]
Como $K_{\alpha_0}$ es compacto, existe una subcubierta finita
\[
K_{\alpha_0}\subset K_{\alpha_1}^c\cup\cdots\cup K_{\alpha_m}^c.
\]
Por lo tanto,
\[
K_{\alpha_0}\cap K_{\alpha_1}\cap\cdots\cap K_{\alpha_m}=\varnothing,
\]
contradicción.
\end{proof}

\begin{corollary}[Intersección anidada]
Sea $\{K_n\}_{n=1}^\infty$ una familia de conjuntos compactos no vacíos tal que
\[
K_1\supset K_2\supset K_3\supset\cdots.
\]
Entonces
\[
\bigcap_{n=1}^\infty K_n\neq\varnothing.
\]
\end{corollary}

\subsection*{Métricas equivalentes}

\begin{definition}[Métricas equivalentes]
Sean $d_1$ y $d_2$ dos métricas sobre un mismo conjunto $X$.
Decimos que $d_1$ y $d_2$ son equivalentes si determinan la misma familia de conjuntos abiertos.
\end{definition}

\begin{example}
En $\mathbb R^k$, sea $d_1$ la métrica euclidiana y sea
\[
d_2(x,y)=\max\{|x_i-y_i|:i=1,\dots,k\}.
\]
Entonces $d_2$ es una métrica.
Para $k=2$,
\[
B^{d_1}_1(0)=\{(x,y)\in\mathbb R^2:\sqrt{x^2+y^2}<1\},
\]
y
\[
B^{d_2}_1(0)=\{(x,y)\in\mathbb R^2:\max\{|x|,|y|\}<1\}.
\]
Estas métricas son equivalentes.
\end{example}

\section*{21May26}

\subsection*{Compactitud en $\mathbb R$}

\begin{theorem}
Si $a,b\in\mathbb R$ con $a<b$, entonces $[a,b]$ es compacto.
\end{theorem}

\begin{proof}
Sea $\mathcal U$ una cubierta abierta de $[a,b]$.
Supongamos que no existe subcubierta finita.
Sea
\[
I_0=[a,b].
\]
Dividimos $I_0$ en dos mitades y elegimos una de ellas que no admita subcubierta finita; la llamamos $I_1$.
Repetimos el procedimiento y obtenemos una sucesión decreciente de intervalos cerrados
\[
I_0\supset I_1\supset I_2\supset\cdots
\]
tal que $\operatorname{long}(I_n)=\dfrac{b-a}{2^n}$ y ninguno de ellos admite subcubierta finita.

Por el teorema de los intervalos encajados, existe
\[
x\in \bigcap_{n=0}^\infty I_n.
\]
Como $\mathcal U$ cubre a $[a,b]$, existe $U\in\mathcal U$ tal que $x\in U$.
Al ser $U$ abierto, existe $\varepsilon>0$ con
\[
(x-\varepsilon,x+\varepsilon)\subset U.
\]
Tomamos $n_0$ tal que
\[
\frac{b-a}{2^{n_0}}<\varepsilon.
\]
Como $x\in I_{n_0}$ y la longitud de $I_{n_0}$ es menor que $\varepsilon$, se tiene
\[
I_{n_0}\subset (x-\varepsilon,x+\varepsilon)\subset U.
\]
Esto da una subcubierta finita de $I_{n_0}$, contradicción.
Luego $[a,b]$ es compacto.
\end{proof}

\begin{corollary}
Un subconjunto $K\subset\mathbb R$ es compacto si y solo si es cerrado y acotado.
\end{corollary}

\begin{proof}
Si $K$ es compacto, entonces es cerrado. Además, la familia
\[
\{(x-1,x+1)\mid x\in K\}
\]
cubre a $K$, y por compacidad se extrae una subcubierta finita; de ahí se deduce que $K$ es acotado.

Recíprocamente, si $K$ es cerrado y acotado, entonces existe $a<b$ tal que $K\subset[a,b]$.
Como $[a,b]$ es compacto y $K$ es cerrado en $\mathbb R$, resulta que $K$ es compacto.
\end{proof}

\subsection*{Métricas equivalentes}

\begin{proposition}[Norma del supremo y norma euclidiana en $\mathbb R^n$]
Sea $n\in\mathbb N$ fijo, y sean $x,y\in\mathbb R^n$.
Entonces
\[
\max_{1\le i\le n}|x_i-y_i|
\le
\left(\sum_{i=1}^n(x_i-y_i)^2\right)^{1/2}
\le
\sqrt n\,\max_{1\le i\le n}|x_i-y_i|.
\]
\end{proposition}

\begin{proof}
Sea
\[
m=\max_{1\le i\le n}|x_i-y_i|.
\]
Entonces
\[
m^2=\max_{1\le i\le n}|x_i-y_i|^2
\le \sum_{i=1}^n(x_i-y_i)^2
\le n\,m^2.
\]
Tomando raíz cuadrada se obtiene la desigualdad.
\end{proof}

\begin{theorem}[Criterio de equivalencia]
Sea $(X,d_1)$ un espacio métrico y sea $d_2$ otra métrica sobre $X$.
Si existen constantes $c,C>0$ tales que
\[
c\,d_1(x,y)\le d_2(x,y)\le C\,d_1(x,y),\qquad \forall x,y\in X,
\]
entonces $d_1$ y $d_2$ son equivalentes.
\end{theorem}

\begin{proof}
Sea $A\subset X$ abierto en $(X,d_1)$ y sea $a\in A$.
Entonces existe $r>0$ tal que
\[
N_r^{d_1}(a)\subset A.
\]
Como
\[
N_{cr}^{d_2}(a)\subset N_r^{d_1}(a),
\]
se tiene $N_{cr}^{d_2}(a)\subset A$; luego $A$ es abierto en $(X,d_2)$.

Recíprocamente, si $A$ es abierto en $(X,d_2)$ y $a\in A$, existe $r>0$ tal que
\[
N_r^{d_2}(a)\subset A.
\]
Pero
\[
N_{r/C}^{d_1}(a)\subset N_r^{d_2}(a),
\]
así que $A$ es abierto en $(X,d_1)$.
\end{proof}

\begin{proposition}[Métrica truncada]
Sea $(X,d_1)$ un espacio métrico y definamos
\[
d_2(x,y)=\frac{d_1(x,y)}{1+d_1(x,y)}.
\]
Entonces $d_2$ es una métrica.
\end{proposition}

\begin{proof}
La positividad y la simetría son inmediatas.
Para la desigualdad del triángulo, se observa que la función
\[
t\longmapsto \frac{t}{1+t}
\]
es creciente en $\mathbb R_+$ y que, para $b>0$ fijo,
\[
\frac{b+c}{1+b+c}\le \frac b{1+b}+\frac c{1+c},\qquad c\ge 0.
\]
Aplicando la desigualdad del triángulo para $d_1$ y esta estimación, se obtiene
\[
d_2(x,y)\le d_2(x,z)+d_2(z,y).
\]
\end{proof}

\subsection*{Densidad y separabilidad}

\begin{definition}[Denso]
Sea $(X,d)$ un espacio métrico y sea $D\subset X$.
Decimos que $D$ es denso si
\[
\overline D=X.
\]
\end{definition}

\begin{definition}[Separable]
Un espacio métrico $(X,d)$ es separable si existe $D\subset X$ tal que $D$ es denso y es numerable.
\end{definition}

\begin{example}
Los conjuntos
\[
E_1=\left\{\frac1n:n\in\mathbb N\right\}\cup\{0\},\qquad
E_2=\left\{2+\frac1n:n\in\mathbb N\right\}\cup\{2\},
\]
\[
E_3=\left\{10+\frac1n:n\in\mathbb N\right\}\cup\{10\}
\]
sirven como ejemplos de conjuntos con punto límite aislado visible en la clausura.

Para cada $j\in\mathbb N$, consideremos
\[
F_j=\left\{\frac1j+\frac1n:n=1,2,\dots\right\}\cup\left\{\frac1j\right\},
\qquad
F_j'=\left\{\frac1j\right\}.
\]
La unión numerable
\[
\bigcup_{j=1}^\infty F_j
\]
ejemplifica una construcción de conjunto numerable.
\end{example}

\section*{25May26}

\subsection*{Recordatorio de métricas equivalentes}

\begin{remark}
En las primeras fotos del día se retoma la métrica truncada
\[
d_2(x,y)=\frac{d_1(x,y)}{1+d_1(x,y)},
\]
ya tratada en la sesión anterior, y el criterio de equivalencia de métricas.
\end{remark}

\subsection*{Compacidad de una celda}

\begin{theorem}[Rudin 2.40]
Toda celda-$d$ de la forma
\[
[a_1,b_1]\times\cdots\times[a_d,b_d]\subset\mathbb R^d
\]
es compacta.
\end{theorem}

\begin{proof}
Supongamos que una cubierta abierta de la celda no admite subcubierta finita.
Dividimos la celda en $2^d$ subceldas cerradas y elegimos una de ellas que tampoco admita subcubierta finita.
Iterando el procedimiento obtenemos una sucesión decreciente de celdas cerradas
\[
J_0\supset J_1\supset J_2\supset\cdots
\]
con lados que tienden a $0$.
Por el teorema de los intervalos encajados, existe un punto
\[
x^\ast\in\bigcap_{n=0}^\infty J_n.
\]
Como la cubierta original es abierta, hay un elemento de la cubierta que contiene a $x^\ast$ y, por tanto, contiene a $J_n$ para $n$ suficientemente grande.
Esto contradice la elección de $J_n$.
Luego la celda es compacta.
\end{proof}

\subsection*{Heine-Borel}

\begin{theorem}[Rudin 2.41]
Sea $K\subset\mathbb R^d$.
Entonces son equivalentes:
\begin{enumerate}[label=\alph*)]
    \item $K$ es cerrado y acotado.
    \item $K$ es compacto.
    \item Todo subconjunto infinito de $K$ tiene un punto límite en $K$.
\end{enumerate}
\end{theorem}

\begin{proof}
El hecho $a)\Rightarrow b)$ se obtiene porque todo conjunto cerrado y acotado cabe en una celda-$d$, y esa celda es compacta; por tanto, $K$ es cerrado en un compacto.

El hecho $b)\Rightarrow c)$ es el Teorema 2.37 de Rudin: todo subconjunto infinito de un compacto tiene un punto límite en el compacto.

Falta probar $c)\Rightarrow a)$.
Si $K$ no es acotado, podemos escoger una sucesión $(x_n)\subset K$ tal que
\[
\|x_n\|>n,\qquad n=1,2,\dots
\]
Entonces $S=\{x_1,x_2,\dots\}$ es infinito, pero no puede tener punto límite: si $y\in\mathbb R^d$ y $r>0$, para $n$ suficientemente grande se tiene
\[
\|x_n-y\|\ge \|x_n\|-\|y\|>r,
\]
de modo que solo finitos términos de $S$ pueden caer en $N_r(y)$.
Contradicción.
Así, $K$ es acotado.

Si $K$ no es cerrado, existe $x_0\in\overline K\setminus K$.
Entonces hay una sucesión $(x_n)\subset K$ con $x_n\to x_0$.
El conjunto $S=\{x_1,x_2,\dots\}$ es infinito y su único punto límite es $x_0\notin K$, contradiciendo $c)$.
Luego $K$ es cerrado.
\end{proof}

\section*{28May26}

\subsection*{Bolas abiertas y cerradas}

\begin{lemma}
Sea $(X,d)$ un espacio métrico. Sean $a\in X$ y $r>0$.
Definimos
\[
V_r(a):=\{x\in X\mid d(x,a)\le r\},
\qquad
N_r(a):=\{x\in X\mid d(x,a)<r\}.
\]
Entonces:
\begin{enumerate}[label=\roman*)]
    \item $V_r(a)$ es cerrado.
    \item $N_r(a)\subset V_r(a)$, aunque la contención puede ser estricta.
    \item Si $X=\mathbb R^k$, entonces $\overline{N_r(a)}=V_r(a)$.
\end{enumerate}
\end{lemma}

\begin{proof}
Para (i), probamos que $(V_r(a))^c$ es abierto.
Sea $b\in (V_r(a))^c$, de modo que $d(a,b)>r$.
Tomamos
\[
d_0=\frac12\bigl(d(a,b)-r\bigr)>0.
\]
Si $x\in N_{d_0}(b)$, entonces
\[
d(a,x)\ge d(a,b)-d(b,x)>r,
\]
de donde $x\notin V_r(a)$.
Así $N_{d_0}(b)\subset (V_r(a))^c$.

La afirmación (ii) es inmediata.

Para (iii), sea $x\in V_r(a)$ y considere el segmento
\[
x_\lambda=(1-\lambda)a+\lambda x,\qquad 0\le\lambda<1.
\]
Entonces
\[
d(a,x_\lambda)=\|a-((1-\lambda)a+\lambda x)\|
=\|\lambda(a-x)\|
=\lambda\|a-x\|
\le \lambda r<r,
\]
por lo que $x_\lambda\in N_r(a)$ para todo $0\le\lambda<1$.
Como $x_\lambda\to x$ cuando $\lambda\to 1$, se concluye que $x\in\overline{N_r(a)}$.
\end{proof}

\begin{remark}
En la métrica discreta puede ocurrir que $N_r(a)\subsetneq V_r(a)$.
\end{remark}

\subsection*{Weierstrass}

\begin{corollary}[Weierstrass]
Todo subconjunto infinito y acotado de $\mathbb R^d$ tiene un punto límite en $\mathbb R^d$.
\end{corollary}

\begin{proof}
Sea $T\subset\mathbb R^d$ infinito y acotado.
Entonces $T$ cabe en una celda-$d$ compacta
\[
[a_1,b_1]\times\cdots\times[a_d,b_d].
\]
Por el Teorema 2.37 de Rudin, todo subconjunto infinito de un compacto tiene un punto límite en el compacto.
Luego $T$ tiene un punto límite en $\mathbb R^d$.
\end{proof}

\subsection*{Normas en $\mathbb R^k$}

\begin{remark}
En el pizarrón aparecen intercambiados los subíndices de dos normas; aquí se normaliza la convención usual.
\end{remark}

\begin{definition}[Norma euclidiana y norma del supremo]
Para $x=(x_1,\dots,x_k)\in\mathbb R^k$ definimos
\[
\|x\|_2=\left(\sum_{i=1}^k x_i^2\right)^{1/2},
\qquad
\|x\|_\infty=\max\{|x_1|,\dots,|x_k|\}.
\]
Para $a\in\mathbb R^k$,
\[
\|x-a\|_2=\left(\sum_{i=1}^k(x_i-a_i)^2\right)^{1/2},
\qquad
\|x-a\|_\infty=\max\{|x_i-a_i|:i=1,\dots,k\}.
\]
\end{definition}

\begin{exercise}
Verificar que la norma del supremo satisface:
\begin{enumerate}[label=\roman*)]
    \item $\|\alpha x\|_\infty=|\alpha|\,\|x\|_\infty$ para todo $x\in\mathbb R^k$ y todo $\alpha\in\mathbb R$.
    \item $\|x\|_\infty=0$ si, y solo si, $x=0$.
    \item $\|x+y\|_\infty\le \|x\|_\infty+\|y\|_\infty$.
\end{enumerate}
\end{exercise}

\end{document}
